\chapter{Implementation and Pseudocode}
\label{ch:impl}

\section{Technology Stack}

The implementation is in Python 3.12. PyTorch and torchvision run the ResNet18 heads; Ultralytics
YOLO11 trains and runs the two detectors; OpenCV and NumPy cut crops and process masks; SciPy supplies
the linear-assignment solver for numbering; pandas joins annotations; scikit-learn computes metrics and
fits the isotonic calibrator; Matplotlib draws the figures; Streamlit serves the demo; pytest runs the
tests; and Docker packages the demo. Exact versions are pinned in \texttt{requirements-lock.txt}
(Table~\ref{tab:hwsw}). NumPy is held below version 2, which breaks the SciPy import that the
Hungarian assignment needs.

\section{Repository Structure}

The repository layout, to two levels, is shown in Figure~\ref{fig:repotree}.

\begin{figure}[!htbp]
\centering
% Repository structure (top levels). Traced to `git ls-files` of the repository on 2 Nov 2026.
\begin{minipage}{0.92\textwidth}
\small
\renewcommand{\DTstyle}{\ttfamily\small}
\renewcommand{\DTcomment}{\rmfamily\itshape\color{p2grey}}
\dirtree{%
.1 An-Intelligent-Dental-Assisting-System/.
.2 app.py\DTcomment{Streamlit web demo}.
.2 Dockerfile\DTcomment{CPU-only container with a build self-test}.
.2 src/\DTcomment{reusable stage modules (orientation, numbering, classifiers, crops)}.
.2 scripts/.
.3 run\_inference.py\DTcomment{Stages 5--6: one OPG in, one JSON report out}.
.3 eval/\DTcomment{40 evaluation scripts}.
.3 dentex/, tufts/, yolo/\DTcomment{dataset builders and trainers}.
.3 figures/\DTcomment{figure scripts with asserted values}.
.3 verify\_paper\_numbers.py\DTcomment{asserts 215 reported numbers}.
.2 models/\DTcomment{one folder per trained stage, with run summaries}.
.2 outputs/\DTcomment{saved evaluation results (JSON, CSV)}.
.2 tests/\DTcomment{23 test modules}.
.2 notebooks/\DTcomment{Phase 2 development notebooks}.
.2 docs/\DTcomment{results log, model metrics, session memory}.
.2 paper/\DTcomment{manuscript and the ten-page IEEE draft}.
.2 requirements.txt, requirements-lock.txt\DTcomment{dependencies}.
}
\end{minipage}

\caption[Repository structure]{Repository structure (top two levels, main folders only).}
\label{fig:repotree}
\end{figure}

The working rules of the repository are part of the implementation. Every result is logged in a
dated, append-only results log. Every default checkpoint is pinned by hash. Superseded checkpoints are
kept, not deleted, so that earlier numbers stay reproducible. Large raw data stay outside version
control.

\section{Modules Carried Over from Phase 2}

\subsection{Module 1: Caries Cascade}

The Phase 2 caries work is a two-stage cascade rather than one three-class model. The first stage, a
ResNet18 on tight tooth crops resized to 224 by 224 with ImageNet normalisation, predicts caries
against no caries; the second reuses the backbone for caries against deep caries, with the final layer
reset and the network fine-tuned at a lower learning rate so that it adapts to severity cues. In
Phase 3 the first stage was retrained with a real negative class of 5,366 healthy teeth, becoming the
retained crop classifier, and the second stage is the shipped severity grader (Stage~3b). % {C084}

\subsection{Module 2: Dataset Exploration and Annotation Parsing}

\texttt{src/data\_loader.py} turns the raw Tufts annotations into a form suitable for tooth-level
analysis. Three parsing bugs were decisive in Phase 2:
\begin{itemize}
\item the expert and box files are lists of records keyed by the stem of their External ID, not
  dictionaries keyed by image name;
\item the boxes are stored as $[y_1, x_1, y_2, x_2]$;
\item the mask values encode lesion instances, not severity classes.
\end{itemize}
The module therefore loads the expert descriptions and derives the image stems, converts the boxes to
$[x, y, w, h]$, classifies each mask as empty or annotated, and joins all sources on the stem.

\subsection{Module 3: Abnormality Dataset Construction}

\texttt{src/crop\_generator.py} and \texttt{src/label\_assigner.py} build the expanded crop from a tooth
box, optionally jitter it, crop the expert mask to the same region and compute the overlap
ratio~\eqref{eq:overlap} that sets the label. This turns sparse region annotations into a tooth-wise
label without lesion segmentation at inference (Figures~\ref{fig:p2grid} and~\ref{fig:p2samples}).

\begin{figure}[!htbp]
\centering
\optionalgraphics[width=0.66\textwidth]{figures/phase2/p2_5_01_crop_validation_grid.jpg}
\caption[Validation grid of tooth boxes, expanded crops and mask-aligned crops]{Validation grid of
tooth boxes, expanded crops and mask-aligned crops (Phase 2 report, Figure 5.1). Tufts radiographs.}
\label{fig:p2grid}
\end{figure}

\begin{figure}[!htbp]
\centering
\optionalgraphics[width=0.86\textwidth]{figures/phase2/p2_5_02_crop_samples.jpg}
\caption[Representative tooth crops of the abnormality dataset]{Representative tooth crops of the
abnormality dataset (Phase 2 report, Figure 5.2). Tufts radiographs.}
\label{fig:p2samples}
\end{figure}

One limitation noted in Phase 2 remains true: the abnormal pool contained jittered near-duplicates of
the same teeth, and offline augmentation multiplied them. The Phase 3 retrain replaced offline
augmentation with an epoch that draws every abnormal crop once and 3.5 normal crops for each, redrawn
every epoch.

\subsection{Modules 4 and 5: Binary Flag and Severity Training}

The Phase 2 binary model fine-tuned a ResNet18 on the expanded crops with focal loss, Adam,
learning-rate reduction on plateau and early stopping; its best epoch was 22. The Phase 2 severity
model was a three-class ResNet18 trained with focal loss on the positive crops. Phase 3 kept both
recipes and changed the split (Section~\ref{sec:sys-stage4}). % {C085}

\subsection{Module 6: CEJ-Based Estimation}

The CEJ heuristic used CLAHE preprocessing, column-wise CEJ and apex detection, a focused search window
for the bone crest, interquartile-range outlier removal and a confidence score
(Figure~\ref{fig:p2cej}). Its results are in Section~\ref{sec:sys-stage4}. % {C086}

\begin{figure}[!htbp]
\centering
\optionalgraphics[width=0.66\textwidth]{figures/phase2/p2_5_03_cej_overlays.jpg}
\caption[CEJ-estimation overlays on tooth crops]{CEJ-estimation overlays on tooth crops from the
exploratory branch (Phase 2 report, Figure 5.3). Tufts radiographs.}
\label{fig:p2cej}
\end{figure}

\section{Modules Added in Phase 3}

\begin{description}[leftmargin=0pt,style=nextline]
\item[Module 7: Tooth detector.] Dataset builders under \texttt{scripts/yolo/} convert each corpus to
  the YOLO format in the detector's rotated frame, balance the corpora and train with Ultralytics.
\item[Module 8: Numbering head.] \texttt{scripts/dentex/build\_enumeration\_crops.py} builds the
  $1.5\times$ crops and box features from DENTEX and Tufts, splitting by source image before any crop
  exists, and \texttt{train\_enumeration\_head.py} trains the two-headed ResNet18.
\item[Module 9: Lesion detector.] \texttt{build\_dentex\_caries\_boxes.py} converts the 705 DENTEX
  disease images into single-class lesion boxes (650 kept), pooled with CariesXrays, with a live MD5
  check against the scored cohort.
\item[Module 10: Integration.] \texttt{scripts/run\_inference.py} composes the stages
  (Algorithm~\ref{alg:pipeline}) and writes the report.
\item[Module 11: Web demo and container.] \texttt{app.py} renders the report in a browser, and the
  Dockerfile packages it with a build-time self-test.
\item[Module 12: Evaluation suite.] Forty scripts under \texttt{scripts/eval/} reproduce every result:
  lesion agreement on DENTEX, Ningxia and CariesXrays; the vision-language and constant baselines;
  numbering, abstention and calibration; the shortcut test; random-effects pooling; constant
  stability; and the near-duplicate check.
\end{description}

\section{Pseudocode}

Eight algorithms give the pseudocode of the pipeline and of the evaluation: end-to-end inference
(Algorithm~\ref{alg:pipeline}), one-to-one FDI assignment (Algorithm~\ref{alg:fdi}), lesion-to-tooth
assignment (Algorithm~\ref{alg:lesion}), the fit of the image-blind constant
(Algorithm~\ref{alg:constant}), the paired cluster bootstrap (Algorithm~\ref{alg:bootstrap}), the
matched-volume comparison (Algorithm~\ref{alg:matched}), random-effects pooling across sites
(Algorithm~\ref{alg:re}) and the near-duplicate check (Algorithm~\ref{alg:nearDup}).

\begin{algorithm}[!htbp]
\caption{End-to-end inference (\texttt{run\_pipeline})}
\label{alg:pipeline}
\begin{algorithmic}[1]
\Require OPG image $I$; loaded models; caries method $m \in \{\text{lesion-detector}, \text{gate}\}$
\Ensure report $R$
\State $W \gets \textsc{CheckInput}(I)$ \Comment{warnings, never a silent failure}
\State $(B, s) \gets \textsc{DetectUpright}(\text{detector}, I)$ \Comment{rotate 90\textdegree, detect at 1,024 px, conf.\ 0.25, NMS 0.50, rotate back}
\State $(c, q) \gets \textsc{AssignFDI}(\text{head}, I, B)$ \Comment{Algorithm~\ref{alg:fdi}}
\If{$m = \text{lesion-detector}$}
  \State $K \gets \textsc{LesionCaries}(I, B)$ \Comment{Algorithm~\ref{alg:lesion}}
\Else
  \State $K \gets \textsc{GateCaries}(I, B)$ \Comment{tight crops, ResNet18, then severity}
\EndIf
\State $A \gets \textsc{AbnormalityFlag}(I, B)$ \Comment{$1.5\times2.0$ crops, ResNet18}
\State \textbf{assert} $|B| = |c| = |K| = |A|$ \Comment{refuse to fuse misaligned outputs}
\For{each tooth $i$}
  \State $\text{assessable}_i \gets (c_i \neq \varnothing) \wedge (q_i \geq 0.30)$
\EndFor
\State $R \gets \textsc{BuildReport}(B, s, c, q, K, A, W)$; sort records by FDI code
\State attach checkpoint hashes, runtime versions and the limitations of the checkpoints that ran
\State \Return $R$
\end{algorithmic}
\end{algorithm}

\begin{algorithm}[!htbp]
\caption{One-to-one FDI assignment (\texttt{assign\_fdi\_head})}
\label{alg:fdi}
\begin{algorithmic}[1]
\Require image $I$ of width $W$ and height $H$; tooth boxes $b_1,\dots,b_n$; two-headed network $g$
\For{$i = 1$ to $n$}
  \State $x_i \gets \textsc{Resize}(\textsc{Crop}(I, \textsc{Expand}(b_i, 1.5, 1.5)), 224)$
  \State $f_i \gets (c_x/W,\ c_y/H,\ w/W,\ h/H)$ of $b_i$
\EndFor
\State $(\ell^{q}, \ell^{p}) \gets g(x, f)$ \Comment{log-softmax over 4 quadrants and 8 positions}
\State $J_{i,(Q,P)} \gets \ell^{q}_{i,Q} + \ell^{p}_{i,P}$ for all 32 codes \Comment{joint log-probability}
\If{any $J$ is not finite} \State \textbf{raise} an error \EndIf
\State $\pi \gets \textsc{LinearSumAssignment}(-J)$ \Comment{Hungarian: maximise $\sum_i J_{i,\pi(i)}$, one code per tooth}
\For{$i = 1$ to $n$}
  \State $c_i \gets \pi(i)$, or none if more than 32 teeth were detected and $i$ is unmatched
  \State $q_i \gets \exp(J_{i,c_i})$; flag $c_i$ as doubtful if $q_i < 0.30$
\EndFor
\State \Return $(c, q)$
\end{algorithmic}
\end{algorithm}

\begin{algorithm}[!htbp]
\caption{Lesion-to-tooth assignment (\texttt{run\_caries\_lesion\_detector})}
\label{alg:lesion}
\begin{algorithmic}[1]
\Require upright image $I$; tooth boxes $b_1,\dots,b_n$; lesion detector $D$; severity network $h$
\State $\mathcal{L} \gets D(I)$ at 1,024 px, confidence $\geq 0.25$ \Comment{no rotation: trained upright}
\State $\text{best} \gets \{\}$
\For{each lesion box $l$ with score $\sigma$ in $\mathcal{L}$}
  \State $r_i \gets \text{area}(l \cap b_i) / \text{area}(l)$ for every tooth $i$ \Comment{containment}
  \If{$\max_i r_i > 0$}
    \State $i^{*} \gets \arg\max_i r_i$
    \State $\text{best}[i^{*}] \gets \max(\text{best}[i^{*}], \sigma)$ \Comment{one tooth per lesion}
  \EndIf \Comment{a lesion outside every tooth is dropped}
\EndFor
\For{each tooth $i$}
  \If{$i \in \text{best}$}
    \State severity $\gets h(\textsc{TightCrop}(I, b_i))$; finding $\gets$ caries or deep caries
  \Else
    \State finding $\gets$ no caries
  \EndIf
\EndFor
\State \Return one finding per tooth, with the detection score marked as such
\end{algorithmic}
\end{algorithm}

\begin{algorithm}[!htbp]
\caption{Fitting the image-blind constant (rule~\eqref{eq:rule})}
\label{alg:constant}
\begin{algorithmic}[1]
\Require fit images, $N$ of them, with carious codes; $n_c$ = images with code $c$ carious; $P = \sum_c n_c$
\State order the 32 codes by $n_c$, largest first
\State $S \gets \varnothing$; $A \gets 0$
\For{each code $c$ in that order}
  \State $k \gets |S|$
  \If{$n_c > 0$ \textbf{and} $2(A + n_c)(kN + P) > 2A((k+1)N + P)$} \Comment{adding $c$ raises F1}
    \State $S \gets S \cup \{c\}$; $A \gets A + n_c$
  \Else
    \State \textbf{break} \Comment{single peak: every later code fails too}
  \EndIf
\EndFor
\State \Return $S$, frozen before it is scored
\end{algorithmic}
\end{algorithm}

\begin{algorithm}[!htbp]
\caption{Paired cluster bootstrap}
\label{alg:bootstrap}
\begin{algorithmic}[1]
\Require images $1,\dots,N$; per-image counts for systems $X$ and $Y$; metric $M$; $R = 10{,}000$; seed 42
\For{$r = 1$ to $R$}
  \State draw $N$ image indices with replacement \Comment{resample radiographs, not teeth}
  \State $\Delta_r \gets M_X(\text{draw}) - M_Y(\text{draw})$ \Comment{both systems on the same draw}
\EndFor
\State \Return point estimate on all images; interval $[\Delta_{(2.5\%)}, \Delta_{(97.5\%)}]$
\end{algorithmic}
\end{algorithm}

\begin{algorithm}[!htbp]
\caption{Matched-volume image-blind comparison}
\label{alg:matched}
\begin{algorithmic}[1]
\Require system predictions on $N$ images; fit-set prevalence order of codes; $P$ carious teeth
\State $v \gets (\text{codes the system names}) / N$; $k \gets \lfloor v \rfloor$; $f \gets v - k$
\For{each image $j$}
  \State $\mathit{TP}^{\text{mix}}_j \gets (1-f)\,\mathit{TP}_j(S_k) + f\,\mathit{TP}_j(S_{k+1})$ \Comment{mixture chosen by image ID, not content}
\EndFor
\State $F_1^{\text{mix}} \gets 2\sum_j \mathit{TP}^{\text{mix}}_j / (vN + P)$
\State $\Delta \gets F_1^{\text{system}} - F_1^{\text{mix}}$; interval by Algorithm~\ref{alg:bootstrap} with $f$ fixed
\State \Return $\Delta$ and its interval
\end{algorithmic}
\end{algorithm}

\begin{algorithm}[!htbp]
\caption{Random-effects pooling across sites, with HKSJ and Q-profile intervals}
\label{alg:re}
\begin{algorithmic}[1]
\Require per-site differences $\Delta_j$ and 95\% intervals $[l_j, u_j]$, $j = 1,\dots,k$
\State $v_j \gets ((u_j - l_j)/3.92)^2$; $w_j \gets 1/v_j$
\State $Q \gets \sum_j w_j (\Delta_j - \bar\Delta_w)^2$; $\hat\tau^2 \gets \max\!\big(0, \frac{Q-(k-1)}{\sum w_j - \sum w_j^2/\sum w_j}\big)$ \Comment{DerSimonian--Laird}
\State $w^{*}_j \gets 1/(v_j + \hat\tau^2)$; $\hat\mu \gets \sum w^{*}_j \Delta_j / \sum w^{*}_j$
\State DL interval: $\hat\mu \pm 1.96\,(\sum w^{*}_j)^{-1/2}$
\State $q \gets \sum w^{*}_j(\Delta_j - \hat\mu)^2/(k-1)$; HKSJ interval: $\hat\mu \pm t_{k-1}\sqrt{q/\sum w^{*}_j}$
\State Q-profile: solve $Q(\tau^2) = \chi^2_{k-1}$ quantiles for the bounds of $\tau^2$
\State prediction interval: $\hat\mu \pm t_{k-2}\sqrt{\hat\tau^2 + \text{SE}^2}$
\State \Return $\hat\mu$, DL and HKSJ intervals, $\tau$ interval, $I^2$, prediction interval
\end{algorithmic}
\end{algorithm}

\begin{algorithm}[!htbp]
\caption{Near-duplicate check (\texttt{dentex\_near\_duplicates.py})}
\label{alg:nearDup}
\begin{algorithmic}[1]
\Require query images (scored cohort); reference images (all other DENTEX images)
\For{each image $x$}
  \State pHash$(x)$: 32$\times$32 greyscale, 2-D DCT, top-left 8$\times$8 thresholded at its median (64 bits)
  \State dHash$(x)$: 9$\times$8 greyscale, sign of the horizontal gradient (64 bits)
\EndFor
\State Hamming distance from every query to every reference, for both hashes
\State positive control: every MD5-identical pair must be at distance 0 under both hashes
\State candidates $\gets$ queries with pHash distance $\leq 6$ or dHash distance $\leq 4$
\State inspect each candidate beside its nearest neighbours by eye; record the verdict
\end{algorithmic}
\end{algorithm}

\section{Verification and Testing}
\label{sec:impl-verify}

The code and the reported numbers were checked in six ways, summarised in Table~\ref{tab:verify}. % {C050}

\begin{table}[!htbp]
\centering
\caption[Verification results]{Verification results of the code and of the reported numbers, re-run
on 6 November 2026 unless stated.}
\label{tab:verify}
\small
\begin{tabularx}{\textwidth}{@{}L{4.0cm}YL{3.4cm}@{}}
\toprule
Check & What it verifies & Result\\
\midrule
Unit and integration tests & 23 modules, including an end-to-end pipeline run & 210 passed, 3 skipped\textsuperscript{a}\\
Number verification & 215 manuscript numbers asserted against their saved artefacts & 215 of 215, no drift\\
Number trace (2 Nov) & every number in the ten-page paper draft traced to a file & 390 numbers; 6 untraced, each explained\\
Invariants & twelve properties of the saved outputs, such as the 86-image exclusion list & 12 of 12 hold\\
Citation audit & every reference resolved on Crossref, OpenAlex, Semantic Scholar or arXiv & verified\\
Container build & imports, the Hungarian solver and the NumPy version, at build time & build self-test\\
\bottomrule
\end{tabularx}
\srcnote{\textsuperscript{a}\,The three skipped tests need the raw Tufts radiographs, which are not on
the machine that ran them.}
\end{table}

\section{Deployment}

The demo container is built with \texttt{docker build -t opg-pipeline .} and run with
\texttt{docker run -p 8501:8501 opg-pipeline}; the page opens at \texttt{http://localhost:8501}. The
build installs the CPU wheel of PyTorch and fails if a dependency would break the pipeline. The running
container is health-checked every 30 seconds. Appendix~B gives the full instructions.
